\documentclass[12pt,spanish,a5paper,landscape]{article}
\usepackage[utf8]{inputenc}
\usepackage{babel}
\usepackage{listings}
\usepackage{mathpazo}
\usepackage{courier}
\usepackage{xcolor}
\usepackage{textcomp}
\usepackage{amsmath}
\usepackage{amssymb}
\usepackage{tikz}
\usepackage{geometry}
\usepackage{enumitem}

\newcommand{\onelinerule}{\rule[2.3ex]{0pt}{0pt}}
\newcommand{\twolinerule}{\rule[6.2ex]{0pt}{0pt}}
\newcommand{\respuesta}{\framebox[\textwidth]{\twolinerule}}
\newcommand{\nombre}{%
  \begin{tikzpicture}[xscale=.4,yscale=.7]
    \draw (0, 0) rectangle (22, 1);
  \end{tikzpicture}%
}
%\newcommand{\rol}   {\framebox[0.3\textwidth]{\onelinerule}}
\newcommand{\rol}{%
  \begin{tikzpicture}[xscale=.4,yscale=.7]
    \draw[gray!40] ( 0, 0) grid      ( 9, 1);
    \draw          ( 0, 0) rectangle ( 9, 1);
    \draw          (10, 0) rectangle (11, 1);
    \draw (9 + .2, .5) -- (10 - .2, .5);
  \end{tikzpicture}%
}
\newcommand{\li}{\lstinline}
\providecommand{\pond}[1]{[{\small\textbf{#1\%}}]}

\lstdefinelanguage{py}{%
  classoffset=0,%
    morekeywords={%
      False,class,finally,is,return,None,continue,for,lambda,try,%
      True,def,from,nonlocal,while,and,del,global,not,with,print,%
      as,elif,if,or,yield,assert,else,import,pass,break,except,in,raise},%
    keywordstyle=\color{black!80}\bfseries,%
  classoffset=1,
    morekeywords={int,float,str,abs,len,raw_input,exit,range,min,max,%
      set,dict,tuple,list,bool,complex,round,sum,all,any,zip,map,filter,%
      sorted,reversed,dir,file,frozenset,open,%
      array,zeros,ones,arange,linspace,eye,diag,dot},
    keywordstyle=\color{black!50}\bfseries,%
  classoffset=0,%
  sensitive=true,%
  morecomment=[l]\#,%
  morestring=[b]',%
  morestring=[b]",%
  stringstyle=\em,%
}

\lstdefinelanguage{testcase}{%
  moredelim=[is][\bfseries]{`}{`},%
  backgroundcolor=\color{gray!20},%
}

\lstdefinelanguage{file}{%
  frame=single,%
}

\lstset{language=py}
\lstset{basicstyle=\ttfamily}
\lstset{columns=fixed}
\lstset{upquote=true}
\lstset{showstringspaces=false}
\lstset{rangeprefix=\#\ }
\lstset{includerangemarker=false}

\newlist{certamen}{enumerate}{1}
\setlist[certamen]{%
  label=\arabic*.,
  font=\LARGE\bfseries,%
  labelindent=-.5in,%
  leftmargin=0pt,%
  labelsep=1em%
}


\lstset{language=py}

\parskip 1ex
\parindent 0pt

\begin{document}
  \pagestyle{empty}
  \thispagestyle{empty}

  \part*{Control 3, miércoles 1--2}
  \newpage
  La Asociación de Tsunamología mide el nivel de las mareas
  en varias ciudades del país,
  y estima la probabilidad de que haya un tsunami
  calculando qué porcentaje de esas mediciones
  supera un umbral definido para cada ciudad.

  \begin{enumerate}[
    labelindent=-.5in,%
    leftmargin=0pt,%
    labelsep=1em%
  ]
    \item
      Escriba la función \li!probabilidad_tsunami(mareas, umbral)!,
      cuyos parámetros son la lista de mediciones de mareas
      y el umbral de una ciudad,
      que retorne el porcentaje de probabilidad estimada
      de que haya tsunami en esa ciudad.
      \lstinputlisting[
        linerange=CASO\ 1-FIN\ CASO\ 1,
      ]{miercoles-1-2.py}

    \item
      El alcalde de Valparaíso decidió que se evacuaría la ciudad
      si la probabilidad de tsunami supera el 32\%.
      El umbral de mareas para Valparaíso es 270.0.

      Escriba la función \li!hay_que_evacuar(mareas)!,
      que reciba como parámetro la lista de mediciones de mareas,
      y que retorne un valor booleano
      indicando si se debe o no evacuar Valparaíso.
      \lstinputlisting[
        linerange=CASO\ 2-FIN\ CASO\ 2,
      ]{miercoles-1-2.py}
  \end{enumerate}

  \newpage
  \part*{Control 3, miércoles 3--4}
  \newpage
  En el básquetbol hay tres tipos de anotaciones:
  el tiro libre (\li!'L'!), que vale un punto;
  el doble (\li!'D'!), que vale dos puntos; y
  el triple (\li!'T'!), que vale tres puntos.

  \begin{enumerate}[
    labelindent=-.5in,%
    leftmargin=0pt,%
    labelsep=1em%
  ]

    \item
      Escriba la función \li!calcular_puntaje(anotaciones)!,
      cuyo parámetro es la lista de las anotaciones
      de un equipo durante un partido,
      que retorne el puntaje obtenido por el equipo.
      \lstinputlisting[
        linerange=CASO\ 1-FIN\ CASO\ 1,
      ]{miercoles-3-4.py}

    \item
      Escriba la función \li!ganador_partido(anotaciones_a, anotaciones_b)!,
      que reciba como parámetros las listas de anotaciones
      de los equipos A y B,
      y que retorne el nombre del equipo ganador.
      En caso de empate,
      debe retornar \li!'Empate'!.
      \lstinputlisting[
        linerange=CASO\ 2-FIN\ CASO\ 2,
      ]{miercoles-3-4.py}
  \end{enumerate}

  \newpage
  \part*{Control 3, jueves 1--2}
  \newpage
  El carioca es un juego de naipes típico de Chile.
  Cada jugador recibe una mano de varias cartas,
  numeradas del 1 al 13.

  Cuando un jugador tiene tres cartas con el mismo número,
  se dice que tiene un trío.

  \begin{enumerate}[
    labelindent=-.5in,%
    leftmargin=0pt,%
    labelsep=1em%
  ]

    \item
      Escriba la función \li!contar_valores(cartas)!,
      que reciba como parámetro una lista de cartas,
      y retorne una lista de tamaño 13
      con las cantidades de veces que aparece cada número.
      \lstinputlisting[
        linerange=CASO\ 1-FIN\ CASO\ 1,
      ]{jueves-1-2.py}

    \item
      Escriba la función \li!tiene_trio(cartas)!
      que retorne un valor booleano
      indicando si la lista \li!cartas! tiene o no un trío.
      \lstinputlisting[
        linerange=CASO\ 2-FIN\ CASO\ 2,
      ]{jueves-1-2.py}

  \end{enumerate}

  \newpage
  \part*{Control 3, jueves 3--4}
  \newpage

  Un vendedor sale de Santiago
  para hacer varios viajes.
  Los viajes hacia el norte
  se representan con kilometrajes positivos,
  y hacia el sur con kilometrajes negativos.

  Escriba la función \li!distancia_recorrida(viajes)!,
  que reciba como parámetro la lista de viajes,
  y que retorne la distancia total recorrida por el vendedor.

  Escriba la función \li!llego_de_regreso(viajes)!,
  que reciba como parámetro la lista de viajes,
  y que retorne un valor booleano
  que indique si después de los viajes
  el vendedor llegó o no de regreso al lugar de partida.

  \lstinputlisting[
    linerange=CASO-FIN\ CASO,
  ]{jueves-3-4.py}



\end{document}

